\section{Results}
\label{sec:results}

\begin{table*}[t]
\centering\small
\setlength{\tabcolsep}{3.2pt}
\resizebox{\textwidth}{!}{\begin{tabular}{lrrrrrrrrr}
\toprule
 & \multicolumn{3}{c}{AMI Zh$\rightarrow$En} & \multicolumn{3}{c}{AMI En$\rightarrow$Zh} & \multicolumn{3}{c}{AliMeeting Zh$\rightarrow$En (out-of-domain)} \\
\cmidrule(lr){2-4}\cmidrule(lr){5-7}\cmidrule(lr){8-10}
System & cpASR-BLEU & tcpASR-BLEU & LAAL & cpASR-BLEU & tcpASR-BLEU & LAAL & cpASR-BLEU & tcpASR-BLEU & LAAL \\
\midrule
SeamlessStreaming$^{a}$ & 12.10 & 8.48 & 2.73 & 13.79 & 11.31 & 2.50 & 13.59 & 9.42 & 3.24 \\
X-Translator$^{b}$ & 12.05 & 8.64 & \textbf{2.69} & 23.31 & 19.96 & 2.74 & 19.46 & 13.77 & \textbf{2.77} \\
Cascade$^{c}$ & 24.88 & 18.37 & 3.51 & 35.07 & 30.97 & 2.80 & 21.81 & 16.16 & 3.55 \\
\midrule
B (deadline baseline) & 24.72 & 18.07 & 3.12 & 35.83 & 31.65 & 2.54 & 22.60 & 16.73 & 3.28 \\
\midrule
\sys{} (ours) & \textbf{26.84} & \textbf{20.35} & 3.03 & \textbf{36.49} & \textbf{32.23} & \textbf{2.43} & \textbf{24.04} & \textbf{18.08} & 2.95 \\
\bottomrule
\end{tabular}}
\caption{Multi-speaker speech-to-speech translation of whole test meetings (same online windows and recognisers for all systems). cpASR-BLEU / tcpASR-BLEU ($\uparrow$): speaker-attributed BLEU of the speech without / with time constraints; LAAL ($\downarrow$): speech LAAL in s. Every BLEU difference between \sys{} and another system is significant (Holm-corrected paired test over meetings, $p<0.001$; Table~\ref{tab:syssig}). $^{a}$Public single-speaker SimulS2ST and $^{b}$open speaker-aware cascade, both zero-shot; $^{c}$speaker-blind model trained on our data; $^{a\text{--}c}$attributed by causal relabelling (Appendix~\ref{app:baselines}). Bold: best.}
\label{tab:main}
\end{table*}


\subsection{Comparison with Multi-Speaker S2ST Systems}
\label{sec:res-sys}

Does \sys{} translate meetings better than existing multi-speaker speech-to-speech systems? \sys{} obtains the highest
cpASR-BLEU and tcpASR-BLEU on all three sets, including the out-of-domain AliMeeting meetings (Table~\ref{tab:main}).
The public systems are partly heard earlier (X-Translator on both Zh$\rightarrow$En sets, SeamlessStreaming on AMI)
but keep only a third to three quarters of \sys{}'s tcpASR-BLEU: SeamlessStreaming translates each window as one speaker, and X-Translator tracks speakers only to choose
the cloned voice, so concurrent contributions are lost or misattributed. Their lower
delay is bought with content (Figure~\ref{fig:tradeoff}, Appendix~\ref{app:more}). The same-data cascade, which
translates speaker-blind and attributes afterwards, reaches B's quality but is heard later, whereas \sys{} is both
better and earlier than either on every set; \S\ref{sec:res-overlap}--\ref{sec:res-ablation} trace this to
speaker-first decoding and the playback layer.

\begin{table}[t]
\centering\small
\setlength{\tabcolsep}{2.2pt}
\resizebox{\columnwidth}{!}{\begin{tabular}{lrrrrrrrrr}
\toprule
 & \multicolumn{3}{c}{AMI Zh$\rightarrow$En} & \multicolumn{3}{c}{AMI En$\rightarrow$Zh} & \multicolumn{3}{c}{AliMeeting} \\
\cmidrule(lr){2-4}\cmidrule(lr){5-7}\cmidrule(lr){8-10}
Thinker & all & ov. & CR & all & ov. & CR & all & ov. & CR \\
\midrule
Speaker-blind + relabel & 19.33 & 9.06 & 54.4 & 35.22 & 27.17 & 65.5 & \textbf{12.67} & 8.55 & 63.0 \\
With speaker tokens & \textbf{24.53} & \textbf{17.73} & \textbf{80.2} & \textbf{40.09} & \textbf{35.71} & \textbf{82.3} & 11.47 & \textbf{9.18} & \textbf{80.2} \\
\bottomrule
\end{tabular}}
\caption{Speaker tokens under overlap: text tcpBLEU over all and overlapped (ov.) units and concurrent-speaker recall (CR, \%), Thinkers trained alike with and without (zeroed) speaker tokens. Confirmatory tests (Table~\ref{tab:reg8}): ov.\ and CR gains hold on both AMI sets, only CR on AliMeeting (machine-translated references here; Appendix~\ref{app:data}).}
\label{tab:overlap}
\end{table}


\subsection{Speaker Tokens under Overlapping Speech}
\label{sec:res-overlap}

Do the speaker tokens help, or would attributing a speaker-blind translation afterwards suffice? For each direction,
we train two simultaneous Thinkers alike, with speaker tokens and with them zeroed, and attribute the latter's output
causally with the same diarizer (Table~\ref{tab:overlap}). On both AMI sets, speaker tokens raise tcpBLEU on overlapped
units by about eight and a half points, more than on all units, and concurrent-speaker recall by 17--26 points. This is
the failure of Figure~\ref{fig:motivation}(1): relabelling gives each translated segment one speaker, so once a
speaker-blind model has merged two concurrent contributions no relabelling can split them, whereas speaker-first
decoding (Eq.~\eqref{eq:serial}) translates them as two segments from the start. On the out-of-domain AliMeeting
meetings the transfer is partial: recall still rises by 17 points, but overlapped units are level and all units
slightly lower, so the speaker stream carries over mainly as better separation of concurrent speakers.

\begin{table}[t]
\centering\small
\setlength{\tabcolsep}{2.2pt}
\resizebox{\columnwidth}{!}{\begin{tabular}{llrrrrrrrrr}
\toprule
 & & \multicolumn{5}{c}{EndOffset and its parts (s)} & \multicolumn{4}{c}{speech} \\
\cmidrule(lr){3-7}\cmidrule(lr){8-11}
Set & Sys. & lag & det. & queue & last & total & load & UTMOS & SIM & pauses \\
\midrule
Z$\rightarrow$E & B & 0.67 & 0.66 & 0.65 & 2.10 & 4.08 & 1.37 & 1.80 & 0.827 & 3.3 \\
 & \sys{} & 0.68 & 0.66 & 0.14 & 1.64 & 3.13 & 0.94 & 1.94 & 0.827 & 8.3 \\
\midrule
E$\rightarrow$Z & B & 0.69 & 0.64 & 0.54 & 1.61 & 3.48 & 1.28 & 3.25 & 0.833 & 4.6 \\
 & \sys{} & 0.71 & 0.64 & 0.14 & 1.42 & 2.91 & 0.99 & 3.40 & 0.831 & 9.4 \\
\midrule
Ali & B & 0.88 & 0.40 & 1.33 & 1.96 & 4.57 & 1.56 & 3.48 & 0.873 & 4.4 \\
 & \sys{} & 0.95 & 0.40 & 0.33 & 1.55 & 3.23 & 1.05 & 3.72 & 0.873 & 11.8 \\
\bottomrule
\end{tabular}}
\caption{Where B's delay is removed (test sets, identical text): EndOffset (total) = model lag + detector wait + queue wait (Eq.~\eqref{eq:backlog}) + last block. Load: speech per second of source; UTMOS: automatic MOS predictor; SIM: similarity to the voice prompt; pauses $\ge$0.5\,s per minute.}
\label{tab:decomp}
\end{table}


\subsection{Where the Listener-Side Delay Is Removed}
\label{sec:res-delay}

As \sys{} and B commit identical text, every difference arises in the speech streams. Table~\ref{tab:decomp} splits
EndOffset into the Thinker's lag, the detector's wait, the queue wait and the last block's duration. The first two
barely change. Under B every stream produces more speech than its source holds (load above one) and builds the queue
of Eq.~\eqref{eq:backlog}; the playback layer brings the load to about one, removes most of the queue and shortens the
last block, which together account for the entire reduction of EndOffset. The reduction is largest on AliMeeting,
whose streams queue longest under B: since the queue depends only on speech durations and release times, the playback
layer carries over to out-of-domain meetings without adaptation. The automatic MOS estimate rises and similarity to
each voice prompt is unchanged, so streams keep their speakers' voices; the cost is more frequent pauses between
blocks, which Eq.~\eqref{eq:plpot} does not price.

\begin{table}[t]
\centering\small
\setlength{\tabcolsep}{2.4pt}
\resizebox{\columnwidth}{!}{\begin{tabular}{lrrrrrr}
\toprule
 & \multicolumn{3}{c}{tcpASR-BLEU $\uparrow$} & \multicolumn{3}{c}{EndOffset (s) $\downarrow$} \\
\cmidrule(lr){2-4}\cmidrule(lr){5-7}
Speech layer & Z$\rightarrow$E & E$\rightarrow$Z & Ali & Z$\rightarrow$E & E$\rightarrow$Z & Ali \\
\midrule
R0 B (base Talker) & 18.07 & 31.65 & 16.73 & 4.08 & 3.48 & 4.57 \\
R1 + read merging & 18.32 & 31.67 & 16.62 & 4.00 & 3.45 & 4.41 \\
R2 + backlog-aware rate & 19.88 & 31.91 & 17.49 & 3.31 & 3.05 & 3.47 \\
R3 + PA Talker & 20.17 & 31.93 & 17.82 & 3.13 & 2.90 & 3.24 \\
\midrule
R4 + best of two = \sys{} & \textbf{20.35} & \textbf{32.23} & \textbf{18.08} & \textbf{3.13} & \textbf{2.91} & \textbf{3.23} \\
\bottomrule
\end{tabular}}
\caption{Ablation on the deadline Thinker's text (text tcpBLEU 25.02 / 41.15 / 21.43 in all rows). R2: rate of Eq.~\eqref{eq:rate}; R3: PA Talker; R4: best of two by $f$. \sys{} $-$ B is significant (Table~\ref{tab:robust}); other rows are descriptive.}
\label{tab:ablation}
\end{table}


\subsection{Ablation of the Playback Layer}
\label{sec:res-ablation}

Which component produces these gains (Table~\ref{tab:ablation})? Read-level merging alone changes little. The
backlog-aware rate of Eq.~\eqref{eq:rate} gives the largest step in quality and EndOffset, as it keeps each stream's
load near one and so prevents most of the queue. The PA Talker, trained on the CRF posterior, shortens the speech
further on every set and raises tcpASR-BLEU on both Zh$\rightarrow$En sets, whose streams queue longest, while
En$\rightarrow$Zh quality is unchanged. This matches what the backward messages of Eq.~\eqref{eq:w} reward: a shorter
rendering gains weight where later blocks would otherwise wait. Choosing the more intelligible of two samples adds
quality without changing EndOffset: delay and intelligibility come from different components.

\begin{table}[t]
\centering\small
\setlength{\tabcolsep}{3pt}
\resizebox{\columnwidth}{!}{\begin{tabular}{lrrrr}
\toprule
 & AMI & AMI & \multicolumn{2}{c}{AliMeeting (out-of-domain)} \\
\cmidrule(lr){4-5}
\sys{} $-$ B & Zh$\rightarrow$En & En$\rightarrow$Zh & test & locked \\
\midrule
tcpASR-BLEU, main recognisers & +2.28 & +0.58 & +1.35 & +2.23 \\
tcpASR-BLEU, second recogniser & +1.32 & +3.74 & +0.93 & -- \\
COMET of the speech & +2.58 & +0.34 & +1.47 & -- \\
LAAL (s) & $-$0.09 & $-$0.11 & $-$0.33 & $-$0.37 \\
EndOffset (s) & $-$0.96 & $-$0.58 & $-$1.34 & $-$1.43 \\
\bottomrule
\end{tabular}}
\caption{\sys{} $-$ B with a second recogniser (not used for selection), COMET, and on eight locked AliMeeting meetings never decoded before. All differences are significant (Appendix~\ref{app:more}).}
\label{tab:robust}
\end{table}


\subsection{Robustness and Out-of-Domain Generalisation}
\label{sec:res-robust}

Could the gains over B reflect speech that suits our recognisers, or meetings used during development? A second
recogniser of another family, which is not the CTC model used for selection, and COMET both prefer \sys{} on every set
(Table~\ref{tab:robust}). On eight AliMeeting meetings never decoded before, with systems and tests fixed in advance,
\sys{} again improves tcpASR-BLEU, LAAL and EndOffset, by margins at least as large as on the AliMeeting test set. The
playback layer, trained and set on AMI only, thus carries over to real Mandarin meetings of another corpus.
